\section*{4 Main result: NP-hardness of circuit minimization for multi-output functions}
\subsection*{4.1 Definitions}
\subsubsection*{4.1.1 Multi-output Functions, Concatenations, and Truth Tables}
For a multi-output Boolean function $f:\{0,1\}^n\to\{0,1\}^m$, the components of $f$ are the single-output functions $f_1,\ldots,f_m$ where $f_i:\{0,1\}^n\to\{0,1\}$ is defined so $f_i(x)$ equals the $i$th output bit of $f(x)$.
For a multi-output Boolean function $f$, we let the circuit complexity of $f$, denoted $\CCirc(f)$, be the minimum size of any circuit computing $f$.
For strings $x,y\in\{0,1\}^*$, we let $x\bullet y$ denote the concatenated string. We identify a multi-output Boolean function $f:\{0,1\}^n\to\{0,1\}^m$ with the concatenated string $T_1\bullet\cdots\bullet T_m\in\{0,1\}^{m\cdot2^n}$ where $T_1,\ldots,T_m$ are the truth tables of the components $f_1,\ldots,f_m$ respectively.
If $f:\{0,1\}^n\to\{0,1\}^{m_1}$ and $g:\{0,1\}^n\to\{0,1\}^{m_2}$ are Boolean functions with the same number of inputs, we define the concatenated Boolean function $f\bullet g:\{0,1\}^n\to\{0,1\}^{m_1+m_2}$ given by $f\bullet g(x)=f(x)\bullet g(x)$.
We also use the $\mathop{\bullet}$ symbol to indicate concatenation in a similar way that $\sum$ acts for addition. For example, if $T_1,\ldots,T_m$ are truth tables of functions with the same number of inputs, then we use the notation $\mathop{\bullet}_{i\in[m]}T_i$ to indicate $T_1\bullet\cdots\bullet T_m$.
\subsubsection*{4.1.2 The Evaluation Function and Multi-output Computation}
Each circuit $C$ induces a multi-output Boolean function we call the Evaluation Function of $C$, denoted $\operatorname{Eval}\text{-}C$, that computes the outputs of each of the gates in $C$. In more detail, for a circuit $C$ that takes $n$ inputs and has $s$ gates, the evaluation function induced by $C$, denoted $\operatorname{Eval}\text{-}C:\{0,1\}^n\to\{0,1\}^{s+n}$, is given by $x_1\bullet\cdots\bullet x_n\bullet g_1\bullet\cdots\bullet g_s$ where $x_i$ is the function computed by the $i$th input wire of $C$ and $g_j$ is the function computed by the $j$th gate in $C$ (for this to be well-defined, we need to fix an ordering of the gates of $C$, but, for our purposes, any ordering will do).
Using the Evaluation Function, an equivalent definition of multi-output circuit computation to the one given in the introduction is that a Boolean circuit $C$ computes a (multi-output) Boolean function $f$ if and only if every component of $f$ is a component of $\operatorname{Eval}\text{-}C$.
\subsubsection*{4.1.3 Windows of Truth Tables}
Given a truth table $T$ of length $n$ and a subset $S\subseteq[n]$, we define the $S$-window of $T$, denoted $T\langle S\rangle$, to be the truth table of length $n$ that (informally) ``sees'' $T$ on the elements of $S$ and zeroes everywhere else. Rigorously,
\[T\langle S\rangle(x)=\begin{cases}T(x)&\text{if }x\in S,\\0&\text{otherwise}.\end{cases}\]
\subsubsection*{4.1.4 Canonical DNF Circuits}
For each (single output) Boolean function $f:\{0,1\}^n\to\{0,1\}$, it will be useful to fix an algorithm that outputs a single ``canonical'' Boolean circuit for computing $f$.
For our purposes, many algorithms are possible, but, for concreteness, we will define the canonical circuit of a Boolean function $f:\{0,1\}^n\to\{0,1\}$ to be the naive DNF formula for $f$, denoted $\operatorname{DNF}_f$, given by
\[\operatorname{DNF}_f\stackrel{\mathrm{def}}=((x_1=y_1^1)\land\cdots\land(x_n=y_n^1))\lor\cdots\lor((x_1=y_1^t)\land\cdots\land(x_n=y_n^t))\]
where
\begin{itemize}
\item $y^1,\ldots,y^t$ are the YES inputs of $f$ in lexicographical order,
\item $x_1,\ldots,x_n$ index the bits of the input string $x$,
\item $y_1^j,\ldots,y_n^j$ index the bits of $y^j$ for each $j\in[t]$,
\item and for $i\in[n]$ and $j\in[t]$, $(x_i=y_i^j)$ is syntax for $\begin{cases}x_i&\text{if }y_i^j=1,\\\neg x_i&\text{if }y_i^j=0.\end{cases}$
\end{itemize}
It is easy to see that $\operatorname{DNF}_f$ can be computed in polynomial-time given the truth table of $f$.
Moreover, reading the above definition of $\operatorname{DNF}_f$ from left to right gives a natural way of defining the $k$th gate in $\operatorname{DNF}_f$, whereby the $k$th gate corresponds to the $k$th gate symbol appearing in the above formula. This fact will later be useful in our analysis.
\subsubsection*{4.1.5 Lifting Sets}
Our reduction will use a way to lift subsets into subsets on larger ground sets. To do this, we will first define a canonical partition of $[m]$ into $n$ parts for $m\geq n$. Let $\mathcal P^{m,n}=(P_1^{m,n},\ldots,P_n^{m,n})$ be the partition of $[m]$ into $n$ parts given by
\[P_i^{m,n}=\{j\in[m]:j\equiv i\pmod n\}.\]
From this partition, we can now lift subsets as follows. Let $n\leq m\in\mathbb N$. Let $S\subseteq[n]$. The $m$-lift of $S$, denoted $S^m$, is the set given by
\[S^m\stackrel{\mathrm{def}}=\bigcup_{i\in S}P_i^{m,n}.\]
